\documentclass[11pt,a4paper]{article}

% Pakete
\usepackage[utf8]{inputenc}
\usepackage[T1]{fontenc}
\usepackage[ngerman]{babel}
\usepackage[left=2cm,right=2cm,top=2cm,bottom=2cm]{geometry}
\usepackage{helvet}
\renewcommand{\familydefault}{\sfdefault}
\usepackage{amsmath,amssymb}
\usepackage{graphicx}
\usepackage{tikz}
\usetikzlibrary{arrows.meta,positioning}
\usepackage{tcolorbox}
\tcbuselibrary{skins,breakable}
\usepackage{enumitem}
\usepackage{tabularx}
\usepackage{colortbl}

% Fachfarbe Physik
\definecolor{physik}{HTML}{2c5aa0}
\definecolor{physiklight}{HTML}{e3f2fd}
\definecolor{loesung}{HTML}{c62828}

% Box-Stile
\tcbset{
    merkbox/.style={
        colback=physiklight,
        colframe=physik,
        fonttitle=\bfseries,
        title=#1,
        boxrule=1pt,
        arc=2pt,
        left=5pt,
        right=5pt,
        top=5pt,
        bottom=5pt
    },
    leerbox/.style={
        colback=white,
        colframe=physik,
        boxrule=1pt,
        arc=2pt,
        left=5pt,
        right=5pt,
        top=8pt,
        bottom=8pt,
        height=#1
    }
}

% Kopfzeile
\usepackage{fancyhdr}
\pagestyle{fancy}
\fancyhf{}
\lhead{\small Physik 12 GK}
\chead{\small Photoeffekt -- \textcolor{loesung}{MUSTERLÖSUNG}}
\rhead{\small Name: \underline{\hspace{4cm}}}
\renewcommand{\headrulewidth}{0.4pt}

% Keine Einrückung
\setlength{\parindent}{0pt}
\setlength{\parskip}{6pt}

\begin{document}

% Titel
\begin{center}
{\Large\textbf{Der Photoeffekt -- Einstieg}}\\[0.3em]
{\normalsize Stunde 1: Hallwachs-Versuch und das Rätsel des Lichts}\\[0.3em]
{\large\textcolor{loesung}{\textbf{--- MUSTERLÖSUNG ---}}}
\end{center}

\vspace{0.5em}

% ============================================================
% HISTORISCHER KONTEXT
% ============================================================

\begin{tcolorbox}[merkbox={Historischer Hintergrund}]
\textbf{1887:} Heinrich Hertz entdeckt zufällig, dass UV-Licht die Funkenentladung zwischen Elektroden erleichtert.

\textbf{1888:} Wilhelm Hallwachs untersucht das Phänomen systematisch: Eine negativ geladene Zinkplatte verliert ihre Ladung, wenn sie mit UV-Licht bestrahlt wird.

\textbf{Das Problem:} Das klassische Wellenmodell des Lichts kann dieses Phänomen nicht erklären!
\end{tcolorbox}

\vspace{0.5em}

% ============================================================
% AUFGABE 1: SIMULATION ERKUNDEN
% ============================================================

\textbf{Aufgabe 1: Simulation erkunden} \hfill \textbf{(8 BE)}

Öffne die PhET-Simulation ,,Der photoelektrische Effekt'' und führe folgende Experimente durch:

\begin{enumerate}[label=\alph*)]
\item Stelle die Intensität auf Maximum und wähle \textbf{violettes Licht} (hohe Frequenz).\\
Beobachtung: \textcolor{loesung}{Elektronen werden aus dem Metall herausgelöst, Strom fließt.}

\vspace{0.5em}

\item Wechsle zu \textbf{rotem Licht} (niedrige Frequenz) bei gleicher Intensität.\\
Beobachtung: \textcolor{loesung}{Keine Elektronen, obwohl die Intensität hoch ist.}

\vspace{0.5em}

\item Gehe zurück zu violettem Licht und \textbf{reduziere die Intensität} auf Minimum.\\
Beobachtung: \textcolor{loesung}{Elektronen treten weiterhin aus, aber weniger pro Zeiteinheit.}

\vspace{0.5em}

\item Welche Größe entscheidet, \textbf{ob} Elektronen austreten? \hfill $\square$ Intensität \quad \textcolor{loesung}{$\boxtimes$ Frequenz}
\end{enumerate}

\vspace{0.5em}

% ============================================================
% AUFGABE 2: HYPOTHESEN
% ============================================================

\textbf{Aufgabe 2: Widerspruch zum Wellenmodell} \hfill \textbf{(6 BE)}

Nach dem klassischen \textbf{Wellenmodell} sollte gelten:
\begin{itemize}[nosep]
\item Höhere Intensität $\rightarrow$ mehr Energie $\rightarrow$ Elektronen sollten austreten
\item Auch rotes Licht sollte bei genug Intensität Elektronen auslösen
\end{itemize}

\begin{enumerate}[label=\alph*)]
\item Erkläre in eigenen Worten, warum deine Beobachtungen diesem Modell widersprechen:

\begin{tcolorbox}[leerbox=3.5cm]
\textcolor{loesung}{\begin{itemize}[nosep,leftmargin=*]
\item Rotes Licht löst \textbf{keine} Elektronen aus -- egal wie intensiv
\item Violettes Licht löst Elektronen aus -- selbst bei minimaler Intensität
\item \textbf{Widerspruch:} Wellenmodell sagt ,,mehr Intensität = mehr Energie''
\item \textbf{Beobachtung:} Die Frequenz (Farbe) entscheidet, nicht die Intensität!
\end{itemize}}
\end{tcolorbox}

\item Formuliere eine Hypothese: Wie könnte Licht ,,wirklich'' Energie übertragen?

\begin{tcolorbox}[leerbox=3.5cm]
\textcolor{loesung}{\begin{itemize}[nosep,leftmargin=*]
\item Licht überträgt Energie in \textbf{Portionen} -- sog. \textbf{Quanten} (Einzahl: Quant)
\item Die Energie eines Quants hängt von der \textbf{Frequenz} ab
\item Hohe Frequenz $\rightarrow$ \textbf{energiereiches} Quant $\rightarrow$ löst Elektron aus
\item Niedrige Frequenz $\rightarrow$ \textbf{energiearmes} Quant $\rightarrow$ reicht nicht aus
\item Viele energiearme Quanten können \textbf{nicht} ein energiereiches ersetzen!
\end{itemize}}
\end{tcolorbox}
\end{enumerate}

\vspace{0.5em}

% ============================================================
% KASTEN 1: MERKSATZ
% ============================================================

\begin{tcolorbox}[merkbox={Kasten 1: Merksatz -- Photoeffekt}]
\textbf{Definition:} Der Photoeffekt ist das Herauslösen von \textcolor{loesung}{\textbf{Elektronen}} aus einer

Metalloberfläche durch \textcolor{loesung}{\textbf{Licht (elektromagn. Strahlung)}}.

\vspace{0.5em}

\textbf{Zentrale Beobachtung:}

\begin{itemize}[nosep]
\item Die \textcolor{loesung}{\textbf{Frequenz}} des Lichts entscheidet, ob Elektronen austreten.
\item Die Intensität beeinflusst nur die \textcolor{loesung}{\textbf{Anzahl}} der Elektronen.
\end{itemize}

\vspace{0.5em}

\textbf{Offene Frage:} Warum ist die Energie des Lichts von der Frequenz abhängig?

$\rightarrow$ Antwort: Einstein 1905 (nächste Stunde)
\end{tcolorbox}

\vspace{1em}

% ============================================================
% DIFFERENZIERUNG
% ============================================================

\hrule
\vspace{0.5em}

\textbf{Zusatzaufgabe (Niveau 3):} \hfill \textbf{(4 BE)}

In der Simulation kannst du verschiedene Metalle auswählen (Natrium, Zink, Kupfer, Platin).

\begin{enumerate}[label=\alph*)]
\item Finde heraus, bei welchem Metall die ,,Grenzfrequenz'' am niedrigsten ist: \textcolor{loesung}{\textbf{Natrium}}

\item Stelle eine Vermutung auf, wovon diese Grenzfrequenz abhängt:

\textcolor{loesung}{Die Grenzfrequenz hängt vom Material ab -- genauer: von der Energie, die nötig ist, um ein Elektron aus dem Metall zu lösen (Austrittsarbeit). Natrium hat die niedrigste Austrittsarbeit, daher reicht bereits Licht mit niedrigerer Frequenz aus.}
\end{enumerate}

\vspace{1em}

\begin{center}
\small\textit{Gesamtpunktzahl: 18 BE}
\end{center}

\end{document}
